\clearpage
\onecolumn
\appendices
\renewcommand{\theequation}{\Alph{section}.\arabic{equation}}
\renewcommand{\theHequation}{\Alph{section}.\arabic{equation}}
\setcounter{section}{0}
\setcounter{equation}{0}
\setcounter{table}{0}
\renewcommand*{\theHtable}{appendix.\arabic{table}}

\section{State Encoding and Fixed Normalization}
\label{sec:appendix-state-encoding}

The offloading agent receives the 62-entry observation defined in~\eqref{eq:observation-encoding}.
The first 22 entries describe the task, device, and network conditions, as shown in Table~\ref{tab:state-context-features}.
The remaining 40 entries are arranged in four blocks of ten features, one block for each action slot.
Table~\ref{tab:state-action-features} gives the order of features within each block.
For slot \(a=0,\ldots,3\), feature \(k=1,\ldots,10\) show position \(22+10a+k\) in the observation.
Slot \(a=0\) corresponds to onboard execution.
%
To normalize a value \(x\), the encoder divides it by a positive scale \(b\) and clips the result to the interval \([0,1]\).

\begin{equation}
\mathcal N(x;b)=\operatorname{clip}(x/b,0,1).
\label{eq:appendix-fixed-normalization}
\end{equation}
The encoder uses fixed scales throughout training and evaluation.

The first four context features record the task payload in kilobytes and the values of the legacy \acs{cpu}, memory, and storage demand fields.
The encoder divides all four values by \(M_D=2.5D_{\mathrm{cfg}}\), where \(D_{\mathrm{cfg}}\) is the fixed value of the \texttt{data\_per\_device\_kb} environment.
The payload in the first feature is the size of the individual task.
The scaling applies only to the observation inputs.
Compute-rate demand, available capacity, and battery energy retain their physical units in the admission, execution, and reward calculations.

The encoder scales the handover count by the configured maximum number of devices, \(N_{\max}\).
For the distance features, it uses \(L=\max\{1,\sqrt{X_{\max}^2+Y_{\max}^2}\}\), where \(X_{\max}\) and \(Y_{\max}\) are the upper bounds of the city coordinates in metres.
In the evaluated configuration, \(D_{\mathrm{cfg}}=400\), \(M_D=1000\), \(N_{\max}=200\), and \(X_{\max}=Y_{\max}=10{,}000\)~m, giving \(L\approx14{,}142\)~m.
For feature 17, the mobility observer computes the mean distance of the active devices from their position centroid and rounds it to two decimal places before normalization.

\begin{table}[!htbp]
\centering
\caption{Ordered context features. Indices refer to the first 22 observation positions.}
\label{tab:state-context-features}
\small
\setlength{\tabcolsep}{5pt}
\renewcommand{\arraystretch}{1.12}
\begin{tabular}{@{}r p{0.51\linewidth} p{0.36\linewidth}@{}}
\toprule
Index & Input quantity & Encoded value \\
\midrule
1 & Task payload, numerical kilobyte value & \(\mathcal N(\text{payload};M_D)\) \\
2 & Legacy \acs{cpu} demand, \texttt{cpu\_demand} & \(\mathcal N(\text{legacy CPU demand};M_D)\) \\
3 & Legacy memory demand, \texttt{memory\_demand} & \(\mathcal N(\text{legacy memory demand};M_D)\) \\
4 & Legacy storage demand, \texttt{storage\_demand} & \(\mathcal N(\text{legacy storage demand};M_D)\) \\
5 & Service priority (\(p_i\)): 1 for \acs{hrllc}, 2 for \acs{mc}, 3 for \acs{ic} & \(\mathcal N(p_i;3)\) \\
6 & Application deadline, in milliseconds & \(\mathcal N(\text{deadline};1000)\) \\
7 & Device-to-serving-\acs{bs} distance (\(D_i^{\mathrm{BS}}\)), in metres & \(\mathcal N(D_i^{\mathrm{BS}};L)\) \\
8 & Received signal level (\(S_i^{\mathrm{rx}}\)), in dBm & \(\mathcal N(|S_i^{\mathrm{rx}}|;120)\) \\
9 & Device has onboard compute capability & Binary indicator \\
10 & Device handover delay, in milliseconds & \(\mathcal N(\text{handover delay};500)\) \\
11 & Device movement speed, in metres per second & \(\mathcal N(\text{speed};30)\) \\
12 & Predicted dwell time, in milliseconds & \(\mathcal N(\text{dwell time};60{,}000)\) \\
13 & Device handover-risk value & \(\mathcal N(\text{handover risk};1)\) \\
14 & Device remaining-battery fraction \(\beta_i\) & Clipped to \([0,1]\) \\
15 & Number of handovers in the current mobility step & \(\mathcal N(\text{handover count};N_{\max})\) \\
16 & Sum of handover delays in the current mobility step, in milliseconds & \(\mathcal N(\text{handover-delay sum};1000)\) \\
17 & Mean distance of active devices from their position centroid, in metres & \(\mathcal N(\text{position dispersion};L)\) \\
18 & Device horizontal coordinate (\(x_i\)), in metres & \(\mathcal N(x_i;X_{\max})\) \\
19 & Device battery below the configured low-battery fraction & \(\mathbf 1\{\beta_i<0.30\}\) \\
20 & Service indicator: \acs{hrllc} & \(\mathbf 1\{p_i=1\}\) \\
21 & Service indicator: \acs{mc} & \(\mathbf 1\{p_i=2\}\) \\
22 & Service indicator: \acs{ic} & \(\mathbf 1\{p_i=3\}\) \\
\bottomrule
\end{tabular}
\end{table}

Each action block describes the conditions the task would face under that execution option.
Current utilization \(u_{ia}\) measures the fraction of capacity in use before the task arrives.
Projected utilization \(\widehat u_{ia}\) estimates the fraction that would be in use after adding the task.
For onboard execution, both values are based on \acs{cpu} utilization.
For a remote option, the encoder uses the highest utilization among \acs{cpu}, memory, and storage at the selected server.
Current and projected headroom give the corresponding fractions of capacity that are free.

For a remote option, \(g_{ia}\) is the congestion percentage reported for the selected server.
The encoder divides this percentage by 100 and clips the result to \([0,1]\).

Compute adequacy \(\kappa_{ia}\) compares the execution option's total compute capacity \(C_{ia,\mathrm{cpu}}\) with the task's required compute rate \(\rho_i\).
Both are measured in cycles per second.
When both values are positive, the encoder divides capacity by demand and clips the result to \([0,1]\).
The score is zero when capacity is zero or negative.
It is one when capacity is positive and demand is zero or negative.

\begin{table}[H]
\centering
\caption{Ordered features repeated for each of the four action slots.}
\label{tab:state-action-features}
\small
\setlength{\tabcolsep}{5pt}
\renewcommand{\arraystretch}{1.12}
\begin{tabular}{@{}r p{0.51\linewidth} p{0.36\linewidth}@{}}
\toprule
Block index & Input quantity & Encoded value \\
\midrule
1 & Physical admission and route-feasibility mask \(m_{ia}\) & Binary indicator \\
2 & Projected utilization & \(\operatorname{clip}(\widehat u_{ia},0,1)\) \\
3 & Resolved-server congestion percentage & \(\mathcal N(g_{ia};100)\) \\
4 & Projected free headroom & \(1-\operatorname{clip}(\widehat u_{ia},0,1)\) \\
5 & Current free headroom & \(\operatorname{clip}(1-u_{ia},0,1)\) \\
6 & Task-relative compute adequacy & \(\kappa_{ia}\) \\
7 & Remaining-battery fraction of the resolved server (\(\beta_{ia}^{S}\)) & Clipped to \([0,1]\) \\
8 & Resolved-server battery below the low-battery fraction & \(\mathbf 1\{\beta_{ia}^{S}<0.30\}\) \\
9 & Action availability & Binary indicator, defined below \\
10 & Action is onboard execution & \(\mathbf 1\{a=0\}\) \\
\bottomrule
\end{tabular}
\end{table}

For onboard execution, the availability flag is one when local computing is enabled and the device supports it.
The feasibility mask indicates whether the device can accept the particular task under the applicable resource and battery checks.
For a remote option, the availability flag is one when the server-selection procedure finds an operational, reachable server that can accept the task.

To compute the server's remaining-battery fraction, the encoder divides the remaining energy by the battery capacity and clips the result to \([0,1]\).
If there is no selected server or its battery capacity is zero or negative, the encoder sets the fraction to one.
Onboard execution uses this default and records zero server congestion.
The device's own battery fraction appears in feature 14 of Table~\ref{tab:state-context-features}.
